\section*{Capítulo \ref{ch:b3:sequence-stratigraphy} --- Estratigrafía de secuencias}

\begin{solution}{exo:b3:sequence-stratigraphy:1}
$A = 0.5 + 0.1 = \qty{0.6}{m/kyr}$; en \qty{20}{kyr}, \qty{12}{m} de
espacio de acomodación. Con el mar descendiendo, $A = -0.3 + 0.1 = \qty{-0.2}{m/kyr}$: un
descenso relativo, frenado por la subsidencia.
\end{solution}

\begin{solution}{exo:b3:sequence-stratigraphy:2}
(a) \omterm{def:b3:sequence-stratigraphy:tract}{Cortejo de nivel descendente} (\omterm{def:b3:sequence-stratigraphy:regression}{regresión forzada}). (b) \omterm{def:b3:sequence-stratigraphy:tract}{Cortejo transgresivo}. (c)
\omterm{def:b3:sequence-stratigraphy:tract}{Cortejo de nivel alto}. (d) \omterm{def:b3:sequence-stratigraphy:tract}{Cortejo de nivel bajo} (\omterm{def:b3:sequence-stratigraphy:regression}{regresión normal} tras el descenso).
\end{solution}

\begin{solution}{exo:b3:sequence-stratigraphy:3}
Cada \omterm{def:b3:sequence-stratigraphy:parasequence}{parasecuencia} termina más somera que la inferior: un apilamiento progradacional,
con la línea de costa desplazándose hacia el mar, como en un \omterm{def:b3:sequence-stratigraphy:tract}{cortejo de nivel alto} o bajo. El
orden opuesto, cada una terminando más fina y más profunda, es retrogradacional: un
\omterm{def:b3:sequence-stratigraphy:tract}{cortejo transgresivo}.
\end{solution}

\begin{solution}{exo:b3:sequence-stratigraphy:4}
Cinco pares por paquete es la razón entre la excentricidad corta (alrededor de
\qty{100}{kyr}) y la precesión (alrededor de \qty{20}{kyr}): cada par es un
ciclo de precesión, cada paquete un ciclo de excentricidad. La tasa es
$1.6/100 = \qty{1.6}{cm/kyr}$, y un par de \qty{0.32}{m} dura alrededor de
\qty{20}{kyr}, lo que es coherente.
\end{solution}

\begin{solution}{exo:b3:sequence-stratigraphy:5}
$\mathrm dx_s/\mathrm dt = (4000 - 0.15\times20\,000)/50 = \qty{20}{m/kyr}$,
hacia el mar. La línea de costa permanece estacionaria cuando $AW = Q$: $A = 4000/20\,000 =
\qty{0.2}{m/kyr}$.
\end{solution}

\begin{solution}{exo:b3:sequence-stratigraphy:6}
$\mathrm dx_s/\mathrm dt = 4000/50 + 0.2/0.01 = 80 + 20 = \qty{100}{m/kyr}$;
el descenso por sí solo explica \qty{20}{m/kyr}, una quinta parte.
\end{solution}

\begin{solution}{exo:b3:sequence-stratigraphy:7}
Retrocede cuando $Q < AW = 0.4\times25\,000 = \qty{10000}{m^3/kyr}$ por
metro. Con $Q = 6000$: $(6000 - 10\,000)/10 = \qty{-400}{m/kyr}$, un retroceso
de \qty{0.4}{km} cada mil años: una anteplaya baja produce una
transgresión rápida.
\end{solution}

\begin{solution}{exo:b3:sequence-stratigraphy:8}
$1/94.9 - 1/123.9 = \qty{0.00247}{\text{ciclos}/kyr}$: un registro más largo que
$1/0.00247 = \qty{406}{kyr}$ los separa, y \qty{600}{kyr} (resolución
0.00167) es suficiente. $1/94.9 - 1/98.9 = 0.00042$: hacen falta alrededor de
\qty{2400}{kyr}, cuatro veces más largo que el registro.
\end{solution}

\begin{solution}{exo:b3:sequence-stratigraphy:9}
$16.4/(20\times40.98) = \qty{0.0200}{m/kyr}$, \qty{2.0}{cm/kyr}; las capas de ceniza
distan $3.1/0.020 = \qty{155}{kyr}$ entre sí.
\end{solution}

\begin{solution}{exo:b3:sequence-stratigraphy:10}
(a) $\tan\alpha = -0.4/200 = -0.002$: descendiendo \ang{0.11}. (b)
$0.2/50 = 0.004$: ascendiendo \ang{0.23}. Una sección sísmica abarca decenas de
kilómetros con una exageración vertical de decenas a centenas, por lo que el ángulo se aprecia;
un acantilado de unos cientos de metros de longitud muestra un ascenso de un metro o menos, perdido en
la variación de las capas.
\end{solution}

\begin{solution}{exo:b3:sequence-stratigraphy:11}
$1\,296\,000/54.725282 = \qty{23.68}{kyr}$ ($g_5$); $/57.924383 =
\qty{22.37}{kyr}$ ($g_2$); $/68.377912 = \qty{18.95}{kyr}$ ($g_4$);
$/67.834313 = \qty{19.11}{kyr}$ ($g_3$). $g_2 - g_5 = 3.199101''$ por año,
\qty{405.1}{kyr}; $g_4 - g_5 = 13.65263''$, \qty{94.9}{kyr}.
\end{solution}

\begin{solution}{exo:b3:sequence-stratigraphy:12}
Ambos generan $0.3\times100 = \qty{30}{m}$ de espacio de acomodación a la misma tasa; con
el mismo aporte la geometría es la misma, ya que el sedimento responde solo al
nivel relativo del mar. Un cambio eustático es igual en todas partes: la misma
secuencia, de la misma edad, en otros márgenes apunta a la primera historia. Una
historia de subsidencia construida a partir de la historia térmica y de carga de la cuenca
(backstripping), o un registro eustático independiente, como el de isótopos de oxígeno,
también los separa.
\end{solution}

\begin{solution}{pb:b3:sequence-stratigraphy:1}
\textbf{1.} $\dot E = 30\times2\pi/400\times\sin(2\pi t/400)$: como máximo
\qty{0.471}{m/kyr}, subiendo en \qty{100}{kyr} y \qty{500}{kyr}, bajando en
\qty{300}{kyr}.

\textbf{2.} \qty{0.020}{m/kyr} en la charnela, 0.037 en \qty{20}{km}, 0.062 en
\qty{50}{km}, \qty{0.120}{m/kyr} en \qty{120}{km}.

\textbf{3.} $A < 0$ cuando $\sin(2\pi t/400) < -0.0346/0.471 = -0.0734$: $t =
200 + (400/2\pi)\arcsin0.0734 = \qty{204.7}{kyr}$.

\textbf{4.} Con $\sigma = 0.0617$: $t = 400 - 63.66\arcsin0.131 =
\qty{391.6}{kyr}$.

\textbf{5.} Unos \qty{187}{kyr} frente a 200: el descenso eustático debe primero
superar a la subsidencia, y deja de imponerse a ella antes del nivel bajo.

\textbf{6.} $0.045\times400 = \qty{18}{m}$, frente a un intervalo eustático de
\qty{60}{m}: un tercio.

\textbf{7.} Hacia tierra desde 50 hasta unos \qty{120}{kyr}, hacia el mar hasta unos
\qty{440}{kyr}, hacia tierra hasta unos \qty{560}{kyr}, hacia el mar de nuevo hasta
\qty{600}{kyr} (los puntos de retorno del modelo son 118.5, 439.5 y
\qty{559}{kyr}; una tabla cada \qty{50}{kyr} solo los acota).

\textbf{8.} $(6000 - 0.30\times15\,000)/40 = \qty{37.5}{m/kyr}$, hacia el mar.

\textbf{9.} $6000/15\,000 = \qty{0.40}{m/kyr}$.

\textbf{10.} $6000/35.7 + 0.424/0.012 = 168.1 + 35.3 = \qty{203}{m/kyr}$,
dentro de \qty{0.5}{\percent} del modelo.

\textbf{11.} El aporte: \qty{83}{\percent} de la velocidad.

\textbf{12.} $(-42.7 - 16.3)/(50.7 - 17.4) = \qty{-1.77}{m/km}$, \ang{0.10}.
Desde su tiempo, el primer punto se ha hundido $0.0345\times395.5 =
\qty{13.6}{m}$, el segundo $0.0623\times208.5 = \qty{13.0}{m}$: casi el
mismo, así que en las cotas actuales el descenso, \qty{59.0}{m}, difiere de la
caída eustática solo en \qty{0.6}{m}. La caída relativa en la línea de costa en
aquel momento, la eustasia menos la subsidencia durante la caída, fue menor.

\textbf{13.} 150 y 200: \omterm{def:b3:sequence-stratigraphy:tract}{cortejo de nivel alto}; 250, 300, 350: \omterm{def:b3:sequence-stratigraphy:tract}{cortejo de nivel descendente}; 400:
\omterm{def:b3:sequence-stratigraphy:tract}{cortejo de nivel bajo}; 450, 500, 550: \omterm{def:b3:sequence-stratigraphy:tract}{cortejo transgresivo} (la línea de costa retrocede justo después de
440); 600: \omterm{def:b3:sequence-stratigraphy:tract}{cortejo de nivel alto}.

\textbf{14.} 204.7 y \qty{391.6}{kyr} frente a 204.5 y 391.5: lo mismo,
dentro del paso del modelo de \qty{0.5}{kyr}.

\textbf{15.} En la plataforma, la superficie expuesta durante la caída: la parte superior de
los depósitos de nivel alto desde la charnela hasta \qty{17}{km}, prolongada sobre los
foresets escalonados del \omterm{def:b3:sequence-stratigraphy:tract}{cortejo de nivel descendente}; en la cuenca, la concordancia correlativa
en la base o en la parte superior del \omterm{def:b3:sequence-stratigraphy:tract}{cortejo de nivel descendente}, por convención. Una
plataforma real estaría cortada por ríos en \omterm{def:b3:sequence-stratigraphy:condensed}{valles incisos}, meteorizada en suelos,
y parcialmente erosionada.

\textbf{16.} $(18.5 - 52.8)/(559 - 439.5) = \qty{-0.29}{km/kyr}$; la
\omterm{def:b3:sequence-stratigraphy:surfaces}{superficie de máxima inundación} se forma al final de ese intervalo.

\textbf{17.} La línea de costa barre rápidamente una llanura costera casi plana:
el aporte queda atrapado cerca de ella, repartido sobre una zona amplia, y cada lugar
recibe poco. Hacia la cuenca, más allá de unos \qty{55}{km} desde el pie del talud, el fondo marino
está desprovisto de aporte: allí se forma una \omterm{def:b3:sequence-stratigraphy:condensed}{sección condensada}.

\textbf{18.} Resolución $1/30 = \qty{0.033}{\text{ciclos}/m}$; Nyquist $1/(2\times
0.01) = \qty{50}{\text{ciclos}/m}$.

\textbf{19.} $7.5/2.0 = 3.75$, 1 y $0.81/2.0 = 0.405$; las razones orbitales
son $405/100 = 4.05$ y $41/100 = 0.41$.

\textbf{20.} \qty{2.0}{m} por \qty{100}{kyr}: \qty{2.0}{cm/kyr}; $0.81/41 =
\qty{1.98}{cm/kyr}$; $7.5/405 = \qty{1.85}{cm/kyr}$. Las longitudes de onda que
el \omterm{def:b3:sequence-stratigraphy:periodogram}{periodograma} puede mostrar cerca de \qty{8}{m} son 10, 7.5 y \qty{6}{m}: el verdadero
\qty{8.1}{m} cae entre dos de ellas.

\textbf{21.} $30/0.020 = \qty{1500}{kyr}$.

\textbf{22.} $23.7\times0.020 = \qty{0.474}{m}$. Muestreado cada
\qty{25}{cm}, la frecuencia de Nyquist es 2 ciclos por metro, una longitud de onda de
\qty{0.5}{m}: la precesión, a 2.11 ciclos por metro, queda por encima de ella y
se pliega hasta $4 - 2.11 = 1.89$, una longitud de onda falsa de \qty{0.53}{m}.

\textbf{23.} Un solo pico encaja con algún periodo y alguna tasa, sea cual sea su
origen; tres picos cuyas razones coinciden con las razones orbitales son independientes de
la tasa, que debe salir entonces igual en cada uno.

\textbf{24.} La línea de costa se desplaza desde \qty{17.37}{km} al inicio de la
caída hasta \qty{52.84}{km} al final de la regresión de nivel bajo: un desplazamiento
máximo hacia la cuenca de \textbf{\qty[text-series-to-math=true]{35.5}{km}}.
\end{solution}
